% Options for packages loaded elsewhere
% Options for packages loaded elsewhere
\PassOptionsToPackage{unicode}{hyperref}
\PassOptionsToPackage{hyphens}{url}
\PassOptionsToPackage{dvipsnames,svgnames,x11names}{xcolor}
%
\documentclass[
  english,
  12pt,
  11pt,
  a4paper,
  openany]{memoir}
\usepackage{xcolor}
\usepackage[top=20mm,left=20mm,bottom=20mm,right=20mm]{geometry}
\usepackage{amsmath,amssymb}
\setcounter{secnumdepth}{5}
\usepackage{iftex}
\ifPDFTeX
  \usepackage[T1]{fontenc}
  \usepackage[utf8]{inputenc}
  \usepackage{textcomp} % provide euro and other symbols
\else % if luatex or xetex
  \usepackage{unicode-math} % this also loads fontspec
  \defaultfontfeatures{Scale=MatchLowercase}
  \defaultfontfeatures[\rmfamily]{Ligatures=TeX,Scale=1}
\fi
\usepackage{lmodern}
\ifPDFTeX\else
  % xetex/luatex font selection
\fi
% Use upquote if available, for straight quotes in verbatim environments
\IfFileExists{upquote.sty}{\usepackage{upquote}}{}
\IfFileExists{microtype.sty}{% use microtype if available
  \usepackage[]{microtype}
  \UseMicrotypeSet[protrusion]{basicmath} % disable protrusion for tt fonts
}{}
\makeatletter
\@ifundefined{KOMAClassName}{% if non-KOMA class
  \IfFileExists{parskip.sty}{%
    \usepackage{parskip}
  }{% else
    \setlength{\parindent}{0pt}
    \setlength{\parskip}{6pt plus 2pt minus 1pt}}
}{% if KOMA class
  \KOMAoptions{parskip=half}}
\makeatother
% Make \paragraph and \subparagraph free-standing
\makeatletter
\ifx\paragraph\undefined\else
  \let\oldparagraph\paragraph
  \renewcommand{\paragraph}{
    \@ifstar
      \xxxParagraphStar
      \xxxParagraphNoStar
  }
  \newcommand{\xxxParagraphStar}[1]{\oldparagraph*{#1}\mbox{}}
  \newcommand{\xxxParagraphNoStar}[1]{\oldparagraph{#1}\mbox{}}
\fi
\ifx\subparagraph\undefined\else
  \let\oldsubparagraph\subparagraph
  \renewcommand{\subparagraph}{
    \@ifstar
      \xxxSubParagraphStar
      \xxxSubParagraphNoStar
  }
  \newcommand{\xxxSubParagraphStar}[1]{\oldsubparagraph*{#1}\mbox{}}
  \newcommand{\xxxSubParagraphNoStar}[1]{\oldsubparagraph{#1}\mbox{}}
\fi
\makeatother


\usepackage{longtable,booktabs,array}
\usepackage{calc} % for calculating minipage widths
% Correct order of tables after \paragraph or \subparagraph
\usepackage{etoolbox}
\makeatletter
\patchcmd\longtable{\par}{\if@noskipsec\mbox{}\fi\par}{}{}
\makeatother
% Allow footnotes in longtable head/foot
\IfFileExists{footnotehyper.sty}{\usepackage{footnotehyper}}{\usepackage{footnote}}
\makesavenoteenv{longtable}
\usepackage{graphicx}
\makeatletter
\newsavebox\pandoc@box
\newcommand*\pandocbounded[1]{% scales image to fit in text height/width
  \sbox\pandoc@box{#1}%
  \Gscale@div\@tempa{\textheight}{\dimexpr\ht\pandoc@box+\dp\pandoc@box\relax}%
  \Gscale@div\@tempb{\linewidth}{\wd\pandoc@box}%
  \ifdim\@tempb\p@<\@tempa\p@\let\@tempa\@tempb\fi% select the smaller of both
  \ifdim\@tempa\p@<\p@\scalebox{\@tempa}{\usebox\pandoc@box}%
  \else\usebox{\pandoc@box}%
  \fi%
}
% Set default figure placement to htbp
\def\fps@figure{htbp}
\makeatother



\ifLuaTeX
\usepackage[bidi=basic]{babel}
\else
\usepackage[bidi=default]{babel}
\fi
% get rid of language-specific shorthands (see #6817):
\let\LanguageShortHands\languageshorthands
\def\languageshorthands#1{}
\ifLuaTeX
  \usepackage[english]{selnolig} % disable illegal ligatures
\fi


\setlength{\emergencystretch}{3em} % prevent overfull lines

\providecommand{\tightlist}{%
  \setlength{\itemsep}{0pt}\setlength{\parskip}{0pt}}



 
\usepackage[]{natbib}
\bibliographystyle{apalike}


\usepackage[english, samf]{ku-frontpage}
\assignment{Seminar paper}
\advisor{Advisor: Jacob Lundbeck Serup}
\let\frontmatter\relax
\let\mainmatter\relax
\let\backmatter\relax
\makeatletter
\@ifpackageloaded{caption}{}{\usepackage{caption}}
\AtBeginDocument{%
\ifdefined\contentsname
  \renewcommand*\contentsname{Table of contents}
\else
  \newcommand\contentsname{Table of contents}
\fi
\ifdefined\listfigurename
  \renewcommand*\listfigurename{List of Figures}
\else
  \newcommand\listfigurename{List of Figures}
\fi
\ifdefined\listtablename
  \renewcommand*\listtablename{List of Tables}
\else
  \newcommand\listtablename{List of Tables}
\fi
\ifdefined\figurename
  \renewcommand*\figurename{Figure}
\else
  \newcommand\figurename{Figure}
\fi
\ifdefined\tablename
  \renewcommand*\tablename{Table}
\else
  \newcommand\tablename{Table}
\fi
}
\@ifpackageloaded{float}{}{\usepackage{float}}
\floatstyle{ruled}
\@ifundefined{c@chapter}{\newfloat{codelisting}{h}{lop}}{\newfloat{codelisting}{h}{lop}[chapter]}
\floatname{codelisting}{Listing}
\newcommand*\listoflistings{\listof{codelisting}{List of Listings}}
\makeatother
\makeatletter
\makeatother
\makeatletter
\@ifpackageloaded{caption}{}{\usepackage{caption}}
\@ifpackageloaded{subcaption}{}{\usepackage{subcaption}}
\makeatother
\usepackage{bookmark}
\IfFileExists{xurl.sty}{\usepackage{xurl}}{} % add URL line breaks if available
\urlstyle{same}
\hypersetup{
  pdftitle={Betting against beta},
  pdfauthor={Tobias Dines Schlünssen \& Conrad Frederik Kromann},
  pdflang={en},
  colorlinks=true,
  linkcolor={blue},
  filecolor={Maroon},
  citecolor={Blue},
  urlcolor={Blue},
  pdfcreator={LaTeX via pandoc}}


\title{Betting against beta}
\usepackage{etoolbox}
\makeatletter
\providecommand{\subtitle}[1]{% add subtitle to \maketitle
  \apptocmd{\@title}{\par {\large #1 \par}}{}{}
}
\makeatother
\subtitle{a study of proposition 1 \& 2 in american equity markets}
\author{Tobias Dines Schlünssen \& Conrad Frederik Kromann}
\date{14 March 2026}
\begin{document}
\frontmatter
\maketitle


\mainmatter

% --- Info page (KU template page 2) ---
\thispagestyle{empty}
\vspace*{2cm}
\begin{tabular}{@{}p{3.5cm} p{9cm}}
  Institute:          & Department of Economics \\[1ex]
  Faculty:            & Faculty of Social Sciences (SAMF) \\[1ex]
  Author(s):          & Tobias Dines Schlünssen (GFS189)\\
                        & Conrad Frederik Kromann (RTL959)\\[1ex]
  Title and subtitle: & Betting against beta\\
                        &- A study of proposition 1 \& 2 in american equity markets \\[1ex]
  Description:        & [A short abstract or description of your thesis goes here.] \\[1ex]
  Advisor:            & Jacob Lundbeck Serup\\[1ex]
  Date:               & June 2nd, 2026 \\
\end{tabular}
\clearpage
% Abstract page
\begin{abstract}
In this paper we will be examining the market anomoly where stocks with higher betas tends to have lower alphas than those with lower betas. This done by testing the proposition 1 and 2 of from the paper:"" by Frazzini and Pedersen (2014)@FRAZZINI20141. Using data from US equity markets we will test proposition 1 by examiningg the relationship between beta and alpha to test for the presence of the market anonmoly. Then, to test for proposition 2, we will construct a BAB-portfolio to examine the performance of a strategy which exploits this market anomaly by going long on low beta stocks and short on high beta stocks. We find___ 
\end{abstract}

\clearpage

\renewcommand{\contentsname}{Table of content}
\tableofcontents
\clearpage

\chapter{Introduction}\label{introduction}

\section{Background}\label{background}

\emph{{[}Person A{]}} --- Describe the background and motivation for the
analysis here.

\section{Research Questions}\label{research-questions}

\begin{itemize}
\tightlist
\item
  Question 1
\item
  Question 2
\item
  Question 3
\end{itemize}

\chapter{Data}\label{data}

\section{Data Description}\label{data-description}

\emph{{[}Person A{]}} --- Describe your dataset here.

\section{Data Cleaning}\label{data-cleaning}

\chapter{Analysis}\label{analysis}

\section{Exploratory Data Analysis}\label{exploratory-data-analysis}

\emph{{[}Person B{]}} --- Perform and describe your exploratory analysis
here.

\section{Statistical Analysis}\label{statistical-analysis}

\chapter{Results}\label{results}

\section{Visualisations}\label{visualisations}

\emph{{[}Person B{]}} --- Present your results and plots here.

\section{Key Findings}\label{key-findings}

Summarise the key findings from the analysis.

\chapter{Conclusion}\label{conclusion}

\section{Conclusion}\label{conclusion-1}

\emph{{[}Person A or B{]}} --- Summarise what you found and what it
means.

\section{Limitations}\label{limitations}

\begin{itemize}
\tightlist
\item
  Limitation 1
\item
  Limitation 2
\end{itemize}

\section{Future Work}\label{future-work}

Suggestions for further analysis or improvements.

\section{References}\label{references}

\renewcommand{\bibsection}{}
\bibliography{references.bib}


\backmatter



\end{document}
